\chapter{{\fontsize{14pt}{21pt}\selectfont\MakeUppercase{Приложение Г. Руководства пользователя и администратора}}}
\label{app:guides}

\section*{Г.1 Руководство пользователя мобильного приложения (слушатель)}

\subsection*{Г.1.1 Назначение и вход}

Мобильное приложение предназначено для слушателя: просмотр назначенных курсов, прохождение уроков, запуск адаптивного теста, просмотр рекомендаций и истории попыток.\cite{flutter_docs,fastapi_docs,react_docs}

\begin{enumerate}
  \item установить приложение, собранное из каталога \texttt{mobile}, либо получить APK от разработчика;
  \item задать базовый URL API через параметр сборки \texttt{API\_BASE}, указывающий на \texttt{/api/v1} серверной части;
  \item на экране входа ввести адрес электронной почты и пароль учётной записи слушателя (например, демонстрационная учётная запись из README репозитория);
  \item при успешной аутентификации открывается главный контейнер навигации с разделами курсов, рекомендаций и профиля.
\end{enumerate}

\subsection*{Г.1.2 Просмотр курсов и уроков}

\begin{enumerate}
  \item в разделе курсов выбрать назначенный курс;
  \item открыть урок и последовательно просматривать страницы контента; при наличии видеоплеера использовать элементы управления воспроизведением;
  \item по завершении страниц система сохраняет прогресс на серверной части.
\end{enumerate}

\subsection*{Г.1.3 Адаптивный тест}

\begin{enumerate}
  \item на экране курса выбрать действие запуска адаптивного теста;
  \item дождаться загрузки первого вопроса; выбрать вариант ответа и подтвердить отправку;
  \item после каждого ответа серверная часть может изменить целевую сложность следующего вопроса;
  \item по завершении набора вопросов отобразится результат, сводка по темам и рекомендации по повторению.
\end{enumerate}

\section*{Г.2 Руководство администратора и преподавателя (веб-панель)}

\subsection*{Г.2.1 Вход и выбор организации}

\begin{enumerate}
  \item открыть URL веб-панели (локально --- порт \texttt{8081} при Docker Compose из репозитория);
  \item выполнить вход учётной записью администратора организации или преподавателя;
  \item при наличии нескольких организаций выбрать текущую организацию в интерфейсе --- контекст передаётся заголовком для серверной части.
\end{enumerate}

\subsection*{Г.2.2 Управление пользователями и группами}

\begin{enumerate}
  \item открыть раздел пользователей;
  \item создать пользователя, назначить роль в организации;
  \item при необходимости создать группу и включить в неё пользователей для массового назначения курсов.
\end{enumerate}

\subsection*{Г.2.3 Учебный контент}

\begin{enumerate}
  \item создать курс, добавить разделы и уроки;
  \item заполнить страницы урока текстом и медиа; загрузить файлы через раздел медиа при необходимости;
  \item создать тест, добавить вопросы с указанием сложности и связей с темами курса.
\end{enumerate}

\subsection*{Г.2.4 Назначение и аналитика}

\begin{enumerate}
  \item назначить курс пользователю или группе;
  \item открыть раздел аналитики для просмотра сводных показателей и отчётов по слушателям.
\end{enumerate}

\section*{Г.3 Контакты поддержки}

Технические вопросы развёртывания адресуются администратору стенда по внутреннему регламенту организации. Параметры демонстрационных учётных записей приведены в файле \texttt{README.md} репозитория программного комплекса.
